\documentclass[11pt]{article}
\usepackage[margin=1in]{geometry}
\usepackage{amsmath,amssymb,amsthm}
\usepackage{booktabs}

\newcommand{\lc}{c_n}
\newcommand{\chl}{\chi^\lambda}
\newtheorem{conj}{Conjecture}
\newtheorem{prop}{Proposition}
\newtheorem{thm}{Theorem}
\title{Disproof of TLMC Conjecture 00000000485\\
\large The proportion of shapes with $|\chl(\lc)|>1$ does not tend to $1$}
\author{TLMC falsification submission}
\date{October 2, 2026}

\begin{document}
\maketitle

\section{The conjecture}

For a partition $\lambda \vdash n$, let $\chl$ be the irreducible character of
the symmetric group $S_n$ attached to $\lambda$, and let $\lc \in S_n$ be a
cycle of length $n$.

\begin{conj}[00000000485]
The proportion of shapes $\lambda \vdash n$ with $|\chl(\lc)| > 1$ tends to
$1$ as $n \to \infty$.
\end{conj}

\section{The attack: n-cycle characters are always in $\{0,\pm1\}$}

Recall the Murnaghan--Nakayama rule: $\chl(\mu)$ is a signed sum over all rim
hooks (border strips) of size $\mu_1$ that can be removed from $\lambda$,
continued recursively. A rim hook is a connected skew shape containing no
$2\times 2$ block, and each contributes the sign $(-1)^{h-1}$, where $h$ is
its number of rows.

Take $\mu = (n)$, i.e.\ the conjugacy class of long cycles $\lc$. Removing a
rim hook of size $n$ from a diagram with exactly $n$ boxes means removing the
\emph{entire diagram}. This is possible if and only if $\lambda$ itself is a
border strip. For a straight (non-skew) shape this means precisely that
$\lambda$ contains no $2\times 2$ block, i.e.\ $\lambda_2 \le 1$, i.e.\
$\lambda$ is a \emph{hook}
\[
\lambda = (n-k,\, 1^{k}) \qquad (0 \le k \le n-1),
\]
in which case the whole diagram has $h = k+1$ rows and
\[
\chl(\lc) = (-1)^{h-1} = (-1)^{k}.
\]
If $\lambda$ is not a hook, no removal is possible and $\chl(\lc) = 0$.

\begin{thm}[classical]
For every $n \ge 1$ and every $\lambda \vdash n$,
\[
\chl(\lc) \in \{-1,\,0,\,+1\}.
\]
Explicitly, $\chl(\lc) = (-1)^k$ if $\lambda = (n-k,1^k)$ is a hook, and
$0$ otherwise.
\end{thm}

\begin{prop}
Conjecture 00000000485 is false. In fact
\[
\#\{\lambda \vdash n \;:\; |\chl(\lc)| > 1\} = 0 \quad\text{for every } n,
\]
so the proportion in question is the constant function $0$, and
$0 \not\to 1$.
\end{prop}

\begin{proof}
Immediate from the theorem: no shape ever satisfies $|\chl(\lc)| > 1$.
\end{proof}

The conjecture is therefore maximally false: the quantity it asserts tends to
$1$ is identically $0$. There are no boundary cases --- the bound
$|\chl(\lc)| \le 1$ is unconditional for all $n$ and all $\lambda \vdash n$.

\section{Explicit counts}

For small $n$, the number $N(n)$ of shapes with $|\chl(\lc)|>1$, out of
$p(n)$ shapes:

\begin{table}[h]
\centering
\begin{tabular}{ccccccccccccccc}
\toprule
$n$ & 1 & 2 & 3 & 4 & 5 & 6 & 7 & 8 & 9 & 10 & 11 & 12 & 13 & 15\\
\midrule
$p(n)$ & 1 & 2 & 3 & 5 & 7 & 11 & 15 & 22 & 30 & 42 & 56 & 77 & 101 & 176\\
$N(n)$ & 0 & 0 & 0 & 0 & 0 & 0 & 0 & 0 & 0 & 0 & 0 & 0 & 0 & 0\\
\bottomrule
\end{tabular}
\caption{The proportion with $|\chl(\lc)|>1$ is $0$ for every $n$.}
\end{table}

For example, at $n = 10$: the hooks $(10), (9,1), \dots, (1^{10})$ give
$\chi = 1, -1, 1, -1, \dots$ (alternating with $k$), and the remaining $32$
shapes all give $0$.

\section{Verification}

\paragraph{Independent recomputation.} A standalone Python program
(\texttt{reproduce.py}) recomputes all n-cycle characters from scratch via
recursive rim-hook removal, and validates the implementation against three
independent facts: column orthogonality
$\sum_\lambda \chl(\mu)^2 = z_\mu$ for every $\mu \vdash n$; the
hook-length formula $f^\lambda = \chl(1^n)$; and the standard-representation
identity $\chi^{(n-1,1)}(\mu) = \#\mathrm{Fix}(\mu) - 1$. All checks pass for
$n \le 15$; the n-cycle column contains only $0, \pm 1$, and $N(n) = 0$
throughout.

\paragraph{Lean 4 machine check.} The \texttt{lean4/} package contains a
computable model of the Murnaghan--Nakayama value at the long cycle (exact
partition enumeration, border-strip test, alternating sign) and
kernel-checked theorems, with \emph{zero axioms} and \emph{zero}
\texttt{sorry}: the enumerations have the correct cardinalities
($(parts\ 10).length = 42$, $(parts\ 12).length = 77$) and every enumerated
part sums to $n$; the number of shapes with $|\chi| > 1$ equals $0$ for
$n = 3,\dots,12$; and \texttt{\#print axioms} reports that no listed theorem
depends on any axiom.

\section{Conclusion}

\[
\boxed{\text{Conjecture 00000000485 is FALSE: } \Pr_{\lambda\vdash n}\bigl[|\chl(\lc)|>1\bigr] \equiv 0 \not\to 1.}
\]

\end{document}
